\section{Definitions and Lean Signatures}
\label{sec:definitions}

The bridge theorem is stated and proved against polymorphic Lean signatures already present in the repository at the time of writing.
No axiom is redefined for M2; the new contribution is a public obstruction interface that factors Sen's argument through bounded semantic witnesses.
The relevant signatures, summarised here for paper-reading purposes, live in the modules listed in Appendix~\ref{sec:repro}.

\begin{definition}[Social welfare function]
\label{def:swf}
A \emph{social welfare function} on a voter type \(\mathtt{Voter}\) and a finite alternative type \(\mathtt{Alt}\) is a map
\(F : \mathtt{Profile\ Voter\ Alt} \to (\mathtt{Alt} \to \mathtt{Alt} \to \mathtt{Prop})\),
i.e.\ a function from preference profiles to a binary social relation on alternatives.
In Lean this is the type abbreviation \(\mathtt{SWF\ Voter\ Alt}\).
The strict part of the social relation is \(\mathtt{strictPart}\ (F\ p)\), with the standard definition \(R\ x\ y \wedge \neg R\ y\ x\).
\end{definition}

\begin{definition}[Universal preference predicate]
\label{def:prefers}
Given a profile \(p : \mathtt{Profile\ Voter\ Alt}\) and a voter \(i\), the predicate
\(\mathtt{prefers\_i}\ p\ i\ x\ y\) holds when voter \(i\)'s ranking strictly prefers \(x\) to \(y\).
\end{definition}

\begin{definition}[Unanimity (semantic \(\mathtt{UN}\))]
\label{def:un}
\(\mathtt{UN}\ F\) holds when for every profile \(p\) and every pair of alternatives \(x,y\), if every voter prefers \(x\) to \(y\) then \((F\ p)\) strictly prefers \(x\) to \(y\).
This is the polymorphic Lean definition; it does not depend on a particular base case.
\end{definition}

\begin{definition}[Decisiveness and minimal liberalism (semantic \(\mathtt{MINLIB}\))]
\label{def:minlib}
A voter \(i\) is \emph{decisive} on an ordered pair \((x,y)\) if for every profile \(p\), whenever \(i\) prefers \(x\) to \(y\) the social relation strictly prefers \(x\) to \(y\), and similarly with \((y,x)\) (\(\mathtt{Decisive.symm}\)).
\(\mathtt{MINLIB}\ F\) holds when there exist two distinct voters \(i,j\) and alternative pairs \((x,y)\) and \((z,w)\) with \(x \neq y\) and \(z \neq w\), such that \(i\) is decisive on \((x,y)\) and \(j\) is decisive on \((z,w)\).
\textbf{Important:} \(\mathtt{MINLIB}\) does \emph{not} guarantee that \(\{x,y,z,w\}\) are four distinct alternatives; the two decisive pairs may overlap.
\end{definition}

\begin{definition}[Semantic obstruction witness]
\label{def:obstruction-witness}
The structure \(\mathtt{SenRightsWitness}\) packages exactly the semantic data extracted from \(\mathtt{MINLIB}\): two distinct rights holders, two decisive ordered pairs, and the distinguished alternative support \(\{x,y,z,w\}\).
The classified structure \(\mathtt{ClassifiedSenRightsWitness}\) adds the raw overlap branch and induces one of the canonical kinds O2, O3, or O4.
\end{definition}

\begin{definition}[Outcome obstruction]
\label{def:outcome-obstruction}
\(\mathtt{SenOutcomeObstruction}\ F\) is the semantic outcome object produced by the bridge.
It has three constructors: a strict conflict on one ordered pair, a 3-cycle at a profile, or a 4-cycle at a profile.
The theorem \(\mathtt{SenOutcomeObstruction.refutes\_socialAcyclic}\) maps any such obstruction to \(\neg\,\mathtt{SocialAcyclic}\ F\).
\end{definition}

\begin{definition}[Acyclicity and \(\mathtt{SocialAcyclic}\)]
\label{def:acyclic}
A binary relation \(S\) on a finite set is \emph{acyclic} when no element has \(\mathtt{Relation.TransGen}\ S\) back to itself.
The predicate \(\mathtt{SocialAcyclic}\ F\) holds when for every profile \(p\), the strict social preference \(\mathtt{strictPart}\ (F\ p)\) is acyclic.
The helper lemmas \(\mathtt{cycle3\_implies\_not\_acyclic}\) and \(\mathtt{cycle4\_implies\_not\_acyclic}\) are polymorphic and are reused by the obstruction soundness layer.
\end{definition}

\paragraph{The generic bridge signatures.}
The Stage 3 bridge exposes two public theorem names:
\begin{quote}
\code{theorem sen\_outcome\_obstruction\_of\_UN\_MINLIB \char`\{Voter : Type u\char`\}\ \char`\{Alt : Type v\char`\}\ [DecidableEq Alt]\ [Fintype Alt]\ (F : SWF Voter Alt)\ (hUN : UN F)\ (hMINLIB : MINLIB F) : SenOutcomeObstruction F}
\end{quote}
\begin{quote}
\code{theorem sen\_impossibility\_from\_obstruction\_basis \char`\{Voter : Type u\char`\}\ \char`\{Alt : Type v\char`\}\ [DecidableEq Alt]\ [Fintype Alt]\ (F : SWF Voter Alt)\ (hUN : UN F)\ (hMINLIB : MINLIB F) : \char`\~SocialAcyclic F}
\end{quote}
There is no \(\mathtt{Fintype\ Voter}\), no \(\mathtt{DecidableEq\ Voter}\), and no explicit lower-bound hypothesis in these signatures.

\paragraph{The legacy compatibility signature.}
The historical theorem remains available as:
\begin{quote}
\code{theorem sen\_impossibility\_lifts \char`\{n m : \texttt{N}\char`\}\ (\_hn : 2 \leq n)\ (\_hm : 4 \leq m)\ (F : SWF (Fin n) (Fin m))\ (hUN : UN F)\ (hMINLIB : MINLIB F) : \char`\~SocialAcyclic F}
\end{quote}
Its proof body now specializes \(\mathtt{sen\_impossibility\_from\_obstruction\_basis}\).
The hypotheses \(\mathtt{\_hn}\) and \(\mathtt{\_hm}\) are retained solely for backward-compatible alignment with the former statement.
We do not claim that \(\mathtt{UN \wedge MINLIB}\) implies \(\mathtt{4 \leq m}\).

\paragraph{Ambient profiles, not literal subinstances.}
The constructed profiles use two distinguished rights holders and at most two ranking types in the ambient voter type.
Unanimity edges still quantify over all voters in that ambient type.
Thus the safe statement is that the contradiction is supported by two distinguished rights holders, at most four distinguished alternatives, and profiles using at most two ranking types.
We do not claim that every general instance contains a literal two-voter Sen24 subinstance.
